%% kapa: exam=1 round=20 year=2009 period=1 form=A subject=economics subject_ko=경제원론 seq=4 number=44
%% source: Q-Net 감정평가사 기출문제 (cst00309) artlSeq=5024178 file=extracted/1008/5024178/2009년 제20회 감정평가사 1차시험 기출문제/1교시.hwp
%% conversion: kordoc 4.24.51 -> tools/gen1.py (auto); review=proofread-sonnet; flags=none
\begin{question}{44}
A는 매년 여름 방학 중 아르바이트를 하는데 이번 여름의 시간당 임금은 작년보다 50\% 높다. 이에 따라 A는 아르바이트 시간을 줄이고 여가를 더 즐기기로 했다. 다음 설명 중 옳은 것은?
\begin{choices}
  \item 여가에 대한 소득효과가 대체효과보다 크다.
  \item 여가에 대한 대체효과가 소득효과보다 크다.
  \item 여가에 대한 소득효과가 대체효과를 정확히 상쇄한다.
  \item A의 노동공급곡선은 수직선이다.
  \item A의 노동공급곡선은 우상향한다.
\end{choices}
\end{question}
\answer{1}
